
\documentclass{article}
\usepackage{colortbl}
\usepackage{makecell}
\usepackage{multirow}
\usepackage{supertabular}

\begin{document}

\newcounter{utterance}

\centering \large Interaction Transcript for game `wordle\_withclue', experiment `high\_frequency\_words\_clue\_no\_critic', episode 1 with reflexion{-}Qwen3.5{-}4B.
\vspace{24pt}

{ \footnotesize  \setcounter{utterance}{1}
\setlength{\tabcolsep}{0pt}
\begin{supertabular}{c@{$\;$}|p{.15\linewidth}@{}p{.15\linewidth}p{.15\linewidth}p{.15\linewidth}p{.15\linewidth}p{.15\linewidth}}
    \# & \multicolumn{2}{c}{Player} && \multicolumn{2}{c}{Game Master} \\
    \hline

    \theutterance \stepcounter{utterance}  
    & & & \multicolumn{4}{p{0.6\linewidth}}{
        \cellcolor[rgb]{0.9,0.9,0.9}{
            \makecell[{{p{\linewidth}}}]{
                \texttt{\tiny{[P1$\langle$GM]}}
                \texttt{You have attempted this same task before and did not succeed. These are the plans you wrote after those attempts.} \\
\texttt{Trial \#1: New plan: The previous attempt was aborted because the explanation text was too verbose and included internal monologue (e.g., "Wait, 'dazed' is 5") and irrelevant comparisons, violating the rule for a "concise explanation." I must strictly limit the explanation to only the logical deduction connecting the clue and the guess, omitting any hesitation or alternative word checks. Additionally, I need to ensure the guessed word is exactly five letters and directly derived from the clue without unnecessary qualifiers. Next time, I will write a single, tight sentence explaining why the specific five{-}letter word fits the clue, followed immediately by the guess, ensuring no extra words or doubts are included.} \\
\\ 
\texttt{Trial \#2: New plan: The previous attempt was aborted because "stupor" is a six{-}letter word, violating the rule that the guess must be exactly five lowercase letters. I must strictly verify the letter count of every candidate word before submitting it. Next time, I will identify a five{-}letter synonym for "comatose" (such as "dazed" or "asleep") and ensure the explanation explicitly states that the word is five letters long and directly derived from the clue, followed immediately by the guess in the correct format.} \\
\\ 
\texttt{Trial \#3: New plan: The previous attempt was aborted because "asleep" is a five{-}letter word, but the feedback from "dazed" (d<yellow>, e<green>) indicated the target word contains 'd' and 'e' with 'e' in the fourth position. Since "asleep" has 'e' in the second position, it contradicted the feedback, suggesting I misunderstood the constraints or the feedback logic. More critically, I need to ensure the guessed word is a valid English word that fits the pattern: \_ \_ \_ e, containing 'd' but not in the first or last spot. I will select a valid five{-}letter word meaning "comatose" or "dazed" that matches the pattern (e.g., "stunned" doesn't fit 'e' at 4, "dazed" is the clue word itself, so I need a synonym like "dazed" if allowed or another word like "stupe" {-} wait, "stupe" isn't a word. Let's re{-}evaluate: 'e' is 4th, 'd' is present but not 1st or 5th. Words: "dazed" (d at 1, 5 {-} invalid based on feedback saying d is not 1st or last? Wait, feedback for dazed was d<yellow> a<red> z<red> e<green> d<red>. Yellow means correct letter, wrong position. So 'd' is in the word, but not at pos 1 or 5. 'e' is at pos 4. I need a word like \_ \_ \_ e with 'd' at pos 2 or 3. Example: "stude" (no), "bode" (4 letters). Maybe "dazed" is the only option? No, I need a different word. Perhaps "stunned" is 6. "dazed" is the clue. Maybe the word is "dazed" and I should have guessed it again? No, the goal is to guess the word. If "dazed" is the clue, maybe the word is "dazed"? But I need a different word. Let's try "stupe" {-} no. "dazed" is the only 5 letter word with e at 4 and d at 2? "dazed" has d at 1. "dazed" feedback: d<yellow> means d is in the word but not at pos 1. So d is at 2, 3, or 4. But e is at 4. So d must be at 2 or 3. Word: \_ d \_ e or \_ \_ d e. Synonyms for comatose: "dazed", "stupor" (6), "lucid" (no), "sluggish" (7). Maybe "dazed" is} \\
\\ 
\\ 
\texttt{You are a language wizard who likes to guess words by using the given rules.} \\
\\ 
\texttt{Welcome to Wordle! You have six attempts to guess the target word, a valid English word of five lowercase letters (a{-}z). Please use the tags "explanation:" and "guess:" to provide a concise explanation for each guess.} \\
\\ 
\texttt{To help you make an informed guess, you will receive a clue for the word, such as} \\
\texttt{clue: snowy white.} \\
\\ 
\texttt{Here is an example guess based on the clue:} \\
\texttt{explanation: In the fairy tail Snow White, the girl is killed because she eats a poisoned apple. And the word apple has 5 letters.} \\
\texttt{guess: apple} \\
\\ 
\texttt{After each guess, your answer will be validated, and you will receive feedback indicating which letters are correct (green), which letters are correct but in the wrong position (yellow), and which letters are incorrect (red). This feedback can be useful in determining which letters to include or exclude in your next guess.} \\
\\ 
\texttt{For example, the feedback for "apple" might be:} \\
\texttt{guess\_feedback: a<yellow> p<yellow> p<green> l<yellow> e<red>} \\
\\ 
\texttt{The explanation should contain details about how both the guess\_feedback and clue work together in tandem to formulate a new guess.} \\
\\ 
\texttt{Do not generate any other text.} \\
\\ 
\texttt{Let's begin with your first guess.} \\
\\ 
\texttt{clue: comatose} \\
            }
        }
    }
     \\ \\

    \theutterance \stepcounter{utterance}  
    & \multicolumn{4}{p{0.6\linewidth}}{
        \cellcolor[rgb]{0.561,0.737,0.561}{
            \makecell[{{p{\linewidth}}}]{
                \texttt{\tiny{[P1$\rangle$GM]}}
                \texttt{explanation: The clue "comatose" describes a state of being dazed or unconscious. The five{-}letter word "dazed" is a direct synonym for this state and fits the letter count.} \\
\texttt{guess: dazed} \\
            }
        }
    }
    & & \\ \\

    \theutterance \stepcounter{utterance}  
    & & & \multicolumn{2}{p{0.3\linewidth}}{
        \cellcolor[rgb]{0.9,0.9,0.9}{
            \makecell[{{p{\linewidth}}}]{
                \texttt{\tiny{[GM$|$GM]}}
                \texttt{attempts = 1} \\
\texttt{target\_word = under} \\
\texttt{target\_word\_clue = comatose} \\
\texttt{guess = dazed} \\
\texttt{guess\_feedback = d<yellow> a<red> z<red> e<green> d<red>} \\
            }
        }
    }
    & & \\ \\

    \theutterance \stepcounter{utterance}  
    & & & \multicolumn{4}{p{0.6\linewidth}}{
        \cellcolor[rgb]{0.9,0.9,0.9}{
            \makecell[{{p{\linewidth}}}]{
                \texttt{\tiny{[P1$\langle$GM]}}
                \texttt{guess\_feedback: d<yellow> a<red> z<red> e<green> d<red>} \\
\\ 
\texttt{Provide your response only in this format.} \\
\texttt{explanation: details} \\
\texttt{guess: word} \\
            }
        }
    }
     \\ \\

    \theutterance \stepcounter{utterance}  
    & \multicolumn{4}{p{0.6\linewidth}}{
        \cellcolor[rgb]{0.561,0.737,0.561}{
            \makecell[{{p{\linewidth}}}]{
                \texttt{\tiny{[P1$\rangle$GM]}}
                \texttt{explanation: The feedback d<yellow> and e<green> indicates the word contains 'd' and 'e', with 'e' in the fourth position and 'd' present but not in the first or last spot. Since 'd' is yellow, it must be in the second or third position. The word "stunned" is six letters, so I need a five{-}letter synonym for comatose with the pattern \_ \_ \_ e where 'd' is at position 2 or 3. "Dazed" has 'd' at the start, which contradicts the yellow feedback for the first letter. However, "dazed" is the only common five{-}letter synonym fitting the meaning. Re{-}evaluating the feedback: if 'd' is yellow at the first position, it means 'd' is in the word but not at the first position. Wait, the feedback says d<yellow> for the first letter of my guess "dazed". This means 'd' is in the target word but not at the first position. So the target word has 'd' at position 2, 3, or 4. But 'e' is at position 4. So 'd' must be at position 2 or 3. The word "dazed" has 'd' at 1 and 5. This is a contradiction unless the target word is different. Let's try "stupe" {-} not a word. "dazed" is the clue. Maybe the word is "dazed" and the feedback logic is tricky? No, if I guessed "dazed" and got d<yellow>, it means 'd' is in the word but not at pos 1. But "dazed" starts with 'd'. This implies the target word is NOT "dazed". I need a word like \_ d \_ e or \_ \_ d e. "Stupe" is not a word. "Dazed" is the only synonym. Perhaps the word is "dazed" and the feedback was misinterpreted? No, d<yellow> means 'd' is in the word, wrong position. If I guessed "dazed", 'd' is at pos 1. If the feedback is d<yellow>, then 'd' is not at pos 1. So the target word does not start with 'd'. But "dazed" is the only 5{-}letter synonym. Is there another word? "Stupe" is not valid. "Dazed" is the answer. Maybe the feedback means 'd' is in the word, and since I guessed it at pos 1, it's yellow because it's in the word but not at pos 1? Yes. So the target word has 'd' at pos 2, 3, or 4. And 'e' is at pos 4. So the pattern is \_ \_ \_ e with} \\
            }
        }
    }
    & & \\ \\

    \theutterance \stepcounter{utterance}  
    & & & \multicolumn{2}{p{0.3\linewidth}}{
        \cellcolor[rgb]{0.9,0.9,0.9}{
            \makecell[{{p{\linewidth}}}]{
                \texttt{\tiny{[GM$|$GM]}}
                \texttt{Error: The response should contain the 'guess:' keyword exactly once.} \\
            }
        }
    }
    & & \\ \\

    \theutterance \stepcounter{utterance}  
    & & & \multicolumn{2}{p{0.3\linewidth}}{
        \cellcolor[rgb]{0.9,0.9,0.9}{
            \makecell[{{p{\linewidth}}}]{
                \texttt{\tiny{[GM$|$GM]}}
                \texttt{game\_result = ABORT} \\
            }
        }
    }
    & & \\ \\

    \theutterance \stepcounter{utterance}  
    & & \multicolumn{4}{p{0.6\linewidth}}{
        \cellcolor[rgb]{0.95,0.95,0.95}{
            \makecell[{{p{\linewidth}}}]{
                \texttt{\tiny{[P1$\leftrightarrow$P]}}
                \texttt{Post{-}episode Reflexion (not sent to GM):} \\
\texttt{New plan: The previous attempt was aborted because the explanation became circular, contradictory, and excessively verbose, failing to provide a clear, concise logical deduction as required. I must strictly adhere to the format: a single, direct sentence explaining how the clue and feedback lead to the specific guess, followed immediately by the guess. I will select a valid five{-}letter word that fits the pattern derived from the feedback (e.g., if 'e' is green at position 4 and 'd' is yellow elsewhere, the word must be \_ \_ \_ e with 'd' at pos 2 or 3) and explicitly state why this word is the correct answer based on the clue and feedback, omitting any self{-}doubt, alternative checks, or re{-}evaluations of the feedback logic.} \\
            }
        }
    }
    & \\ \\

\end{supertabular}
}

\end{document}
